\subsubsection{Molecular Geometries and Active Space Preparation}
The dataset used comprised 28 closed-shell organic molecules from Thiel's excited-state benchmark set \cite{thiels2008benchmark}. The published Cartesian nuclear coordinates were referenced directly and used without further optimisation. 




\paragraph{Hamiltonian and Spin Operators.}
For each active space, the electronic Hamiltonian was constructed in second quantisation as follows:
\ElectronicHamiltonianEquation
where $E_{\mathrm{nuc}}$ is the core nuclear repulsion energy and $p, q, r, s$ index the active spin-orbitals. $h_{pq}$ and $g_{pqrs}$ are the one- and two-electron integral coefficients evaluated in the fixed canonical Hartree-Fock MO basis.

Spin operators were constructed in the same second quantisation basis as follows.

\SpinZEquation
\SpinLadderOperatorsEquation
\SpinSquaredEquation

All fermionic Hamiltonian, spin and excitation operators were mapped to qubit operators using the Jordan-Wigner (JW) transformation.

\paragraph{Ground State Reference.}
As introduced above, approximate ground-state reference wavefunctions were prepared using parameterised Unitary Coupled-Cluster (UCC) ansatz. 
In this study, six ansatz families were used, namely: UCCSD, UCCSDSinglet, UCCGSD, and \(k\)-UpCCGSDSinglet with \(k=1,2,\) and \(3\). The ansätze were characterised by the excitation generators included in the cluster operator $\hat{T}$, summarised in the following table.

\input{tables/methods/ansatz}

For each ansatz, candidate singles and doubles transitions were first enumerated in fermionic form. Only canonical transitions that excite from a lower to higher spin orbital were retained because the reverse transitions would already be included through $\hat{T}^{\dagger}$ before exponentiation. All retained generators were spin conserving: single excitations preserved spin labels, while double excitations preserved the total number of $\alpha$ and $\beta$ spin electrons. 

The non-generalised UCCSD and UCCSDSinglet ansätze were restricted to occupied-to-virtual single and double excitations relative to the reference Hartree-Fock determinant \cite{Anand2022UCCReview}. In contrast, the generalised UCCGSD ansatz allowed generalised orbital transitions between occupied--occupied and virtual--virtual spin orbitals \cite{Lee2018GeneralisedAnsatzOriginal, Diniz2021GeneralisedAnsatzAdaptations}.
To reduce the size of the generalised-double operator pool, mixed-direction double excitations, in which one electron was promoted while the other was relaxed, were excluded. Across the $4e4o$ to $8e8o$ active spaces, this practical restriction reduced the number of independent parameters by 42.2--45.2\% and the circuit depth by 46.4--47.6\%.

The $k$-UpCCGSDSinglet ansatz combined generalised singles with generalised paired doubles. For each paired double excitation, an $\alpha\beta$ electron pair was removed from one spatial orbital and created in another. 
The UCCSDSinglet and $k$-UpCCGSDSinglet constructions imposed spin-adapted parameter sharing between spin-orbital excitations that had the same spatial-orbital labels. 
For $k$-UpCCGSDSinglet, the complete generalised single and paired-double block was repeated $k$ times, with an independent set of variational parameters assigned to each repetition. 

Unless otherwise stated, all UCC operators were implemented using the first-order Suzuki Trotter product approximation given in \cref{eq:first-order-suzuki-trotter}. 

\SuzukiTrotterFirstOrderEquation

At every optimisation step, the circuit energy was evaluated as the noiseless statevector expectation value of the electronic Hamiltonian \cref{eq:electronic-hamiltonian}.
The variational parameters were optimised classically using the L-BFGS-B optimiser through Qibo, to minimise the ground state expectation value. For the $k$-UpCCGSDSinglet ansatz, the parameters associated with all $k$ repeated layers were optimised simultaneously.
Before optimisation, ansatz parameters were initialised based on the following hierarchy: 
\[
\mathrm{MP2}
\rightarrow
1\text{-UpCCGSD}
\rightarrow
2\text{-UpCCGSD}
\rightarrow
3\text{-UpCCGSD}
\rightarrow
\mathrm{UCCSDSinglet}
\rightarrow
\mathrm{UCCSD}
\rightarrow
\mathrm{UCCGSD}.
\]

Parameters associated with excitation generators shared with the preceding ansatz were initialised with the previously optimised amplitudes, while newly introduced generators were initialised to zero. The initial occupied-to-virtual double-excitation parameters were estimated using second-order Moller-Plesset perturbation theory (MP2). For parameter transfer from a \(k\)-UpCCGSDSinglet ansatz to a subsequent ansatz, the optimized amplitudes associated with each generator were summed across the \(k\) layers and used as the initial amplitude for the corresponding generator

Following optimisation, the resulting VQE parameters were stored as a fixed ground-state reference wavefunction. This fixed reference set was used consistently for all subsequent statevector and finite-shot QSE calculations.



\subsubsection{Numerical Analysis and Evaluation Metrics}
\paragraph{Generalised Eigenvalue Problem.}
Due to the non-orthogonality of the QSE basis states, solving for the energy eigenvalues and wavefunctions involves solving a generalised eigenvalue problem $\mathbf{H}\mathbf{c} = E\mathbf{S}\mathbf{c}$. This was done by first diagonalising the Hermitian overlap matrix $\mathbf{S}$. 

\DiagonaliseSEquation
where the columns of $U$ are the eigenvectors of $\mathbf{S}$ and $\Lambda$ is a diagonal matrix containing the eigenvalues of $\mathbf{S}$. Then, we define the whitening factor

\WhitenSEquation

Using $c=Vy$, we get

\TransformToOrdinaryEigenvalueEquation

Numerical instabilities arise when the eigenvalues used are tiny or near-null, because noise perturbations will be greatly amplified by a factor $\frac{1}{\sqrt{\lambda}}$. 
Therefore in practice, the following reduced eigenvalue equation is solved
\ReducedEigenvalueEquation
where $\tilde{\mathbf{V}}$ and $\tilde{\mathbf{V}}^{\dagger}$ denote a reduced form of $\mathbf{V}$ and $\mathbf{V}^{\dagger}$, retaining only a selected subset of its columns. This gives the wavefunction in the QSE operator basis as $\mathbf{c}_k=\tilde{\mathbf{V}}\mathbf{y}_k$.



\paragraph{Spin-expectation Evaluation.}

For the $k^\mathrm{th}$ recovered state, the spin expectation value was calculated as follows:

\SpinExpectationQSEBasisEquation

Using $\mathbf{c}_k=\tilde{\mathbf{V}}\mathbf{y}_k$ and
$\tilde{\mathbf{V}}^{\dagger}\mathbf{S}\tilde{\mathbf{V}}=\mathbf{I}$, this
expression becomes

\SpinExpectationShortenedBasisEquation
